\subsection{What integral neighbor operators do to the same-genus trace}

This is an audit of a proposed method for detecting the same-discriminant
product.  It supplies no auxiliary field and no rational Selmer vanishing.
It uses full orthogonal isometry classes, as in the senary mass calculation.

\begin{proposition}[KT-NEIGHBOR-CONSTANT-LINE-LIMIT]
Let $\mathcal G$ be a genus of positive-definite even lattices of rank six
and determinant $d$.  Fix a prime $p\geq31$ with $p\nmid2d$, and write
\[
 \Lambda=\mathbf Z_p^{\mathcal G/\cong},\quad
 \langle f,g\rangle=\sum_{[L]}\frac{f(L)g(L)}{|O(L)|},\quad
 M=\langle1,1\rangle.
 \tag{NP.1}
\]
For an odd prime $q\nmid d$, let $T_q$ be the unscaled Kneser neighbor
operator on functions.  Put
\[
 N_q=q^4+q^3+(1+\epsilon_q)q^2+q+1,
 \qquad \epsilon_q=\left(\frac{-d}{q}\right).
 \tag{NP.2}
\]
Then $T_q1=N_q1$.  Consequently every polynomial $P$ with
$\mathbf Z_p$ coefficients in finitely many such neighbor operators satisfies
\[
 \langle P((T_q))\,a1,1\rangle
       =aP((N_q))M\qquad(a\in\mathbf Z_p).
 \tag{NP.3}
\]
If $p\mid M$, these integral neighbor operations followed by genus
averaging never give a unit on a constant input.  In particular
$N_p\equiv1\pmod p$ does not imply $M\in\mathbf Z_p^\times$.

More strongly, if $p\mid M$, there is a simultaneous
characteristic-zero eigenline for all these $T_q$, orthogonal to
every spinor character, whose eigenvalues satisfy
\[
                  \lambda(T_q)\equiv N_q\pmod{\mathfrak p}
                         \quad\hbox{for every odd }q\nmid d
 \tag{NP.4}
\]
at a prime of a finite coefficient extension above $p$.
Its $T_p$ eigenvalue is a unit.
\end{proposition}

\begin{proof}
All automorphism orders in (NP.1) are units at $p$.  Otherwise Cauchy's
theorem would give an order-$p$ element of $GL_6(\mathbf Q)$, whose
nontrivial cyclotomic constituent has dimension $p-1>6$.
Thus the displayed pairing is perfect over $\mathbf Z_p$.

The neighbor construction identifies the neighbors of $L$ with the
isotropic lines in $L/qL$.  Their number is (NP.2), by the rank-six
finite-quadric formula.  This proves the constant row sum and (NP.3).
Reversing an isometry from $L'$ onto a neighbor of $L$ gives the balance
identity
\[
 \frac{\#\{q\text{-neighbors of }L\text{ isometric to }L'\}}{|O(L)|}
 =\frac{\#\{q\text{-neighbors of }L'\text{ isometric to }L\}}{|O(L')|}.
 \tag{NP.5}
\]
Hence the operators are self-adjoint.  Their commutativity is the usual
commutativity of the unramified orthogonal Hecke algebras at each prime,
together with the independence of different primes.

Let $S$ be the group of spinor genera in $\mathcal G$ and put $s=|S|$.
It is a finite elementary abelian two-group; every spinor genus has
mass $M/s$.  For each $u\in S$ choose one lattice $L_u$ in that spinor
genus and set
\[
 \Lambda^0_{\mathrm{sp}}
  =\left\{f\in\Lambda:\sum_{[L]\in u}\frac{f(L)}{|O(L)|}=0
                       \text{ for every }u\in S\right\}.
 \tag{NP.6}
\]
The map of component means is onto $\mathbf Z_p^S$, so this submodule
is saturated.  It is the orthogonal complement of all spinor characters.
A $q$-neighbor shifts the spinor genus by a fixed element $\delta_q$.
Together with (NP.5) and the row sum, this shows that $T_q$ preserves
(NP.6): the mean of $T_qf$ on $u$ is $N_q$ times the mean of $f$ on
$u\delta_q$.  If $p\mid M$, the element
\[
 g=1-\frac Ms\sum_{u\in S}|O(L_u)|\,\mathbf1_{L_u}
       \in\Lambda^0_{\mathrm{sp}}
       \quad\hbox{satisfies}\quad g\equiv1\pmod p.
 \tag{NP.7}
\]
The factor $s$ is a unit.  In particular $g$ is primitive, and its
reduction spans a simultaneous eigenline with eigenvalues $N_q$.

For completeness, let $\mathcal A$ be the image of the algebra generated
by all $T_q$ in $\operatorname{End}_{\mathbf Z_p}(\Lambda^0_{\mathrm{sp}})$.
It is finite and torsion-free over $\mathbf Z_p$, since it is a
submodule of a finite free module, and is commutative.  The line just
constructed defines a character $\mathcal A\to\mathbf F_p$.
A maximal ideal defining that character contains a minimal prime that
avoids $p$: the element $p$ is a nonzerodivisor in $\mathcal A$.
The corresponding finite torsion-free domain has fraction field finite
over $\mathbf Q_p$.  Normalizing and completing at a prime above the
chosen maximal ideal gives a characteristic-zero character reducing to
the specified one.  The operators and the component-mean conditions
are defined over $\mathbf Q$, and their weighted pairing over $\mathbf R$
is positive definite.  Thus the commuting operators are simultaneously
diagonalizable over $\mathbf R$ and hence over an algebraic coefficient
field; scalar extension gives the same conclusion over $\overline{\mathbf Q}_p$.
Therefore this character
occurs on an eigenline of $\Lambda^0_{\mathrm{sp}}$ after finite scalar
extension.  This proves (NP.4).  Finally (NP.2) gives $N_p\equiv1$,
which proves the unit assertion.
\end{proof}

\paragraph{Exact application to the retained genus.}
In KT-SENARY-GENUS-TRACE-DETECTOR, take $p=31$ and $D=-79$.
The already checked values are $M=31/720$ and
$a=-\beta_{17}(\chi_{-79})\equiv11\pmod{31}$.
Here $\epsilon_{31}=1$, since $-79\equiv13^2\pmod{31}$, and
\[
                  N_{31}=962\cdot993=955266\equiv1\pmod{31}.
\]
Thus every expression (NP.3) is zero modulo $31$, although the
constant eigenvalue of $T_{31}$ is a unit.  There is also an actual
simultaneous orthogonal Hecke congruence (NP.4) outside all spinor
characters.  No assertion that this eigenline belongs to the compact
unitary source, or to an admissible deformation space, is made.

\paragraph{Primary inputs and remaining theta-lift issue.}
The primary reference is G.~Chenevier, \emph{Statistics for Kneser
$p$-neighbors}, arXiv:2104.06846v2, 3 May 2021: (1.2)--(1.4) on
p.~3; the full local construction and Lemmas 4.1--4.3 on pp.~13--14;
Definition 4.5 and the self-transpose Hecke algebra on p.~15;
the spinor-genus description and complete proof of Lemma 5.6 on
p.~20; the spinor shift in Remark (R2) on p.~21; and (6.1)--(6.2),
Lemma 6.1 and the self-adjointness paragraph on pp.~22--23.
Those precise passages, including the short local and equal-mass
proofs, were read for this audit.  The mass and Bernoulli values are
imported from the already checked senary trace.

A different theta-lift functional can have varying weights on the
classes of a genus.  It is not covered by the constant row-sum argument
unless it is followed by this same genus average.  Its value on the
constant input is then a different weighted sum, whose integrality,
factorization by the specified same-discriminant product, and unit
value must each be proved.  Existence of an integral theta map alone
proves none of those unit assertions.  This proposition rejects the
neighbor-polynomial and spectral-separation shortcuts above; it does
not rule out a suitable further theta correspondence.
